\Needspace{11\baselineskip}

\section{Cognitive Modeling: Course Blueprint (16 weeks)}

\subsection{Design principles}

\begin{itemize}
\tightlist
\item
  \textbf{One modeling pipeline used all term.} Every lab follows the
  same steps: \emph{simulate \ensuremath{\rightarrow} fit \ensuremath{\rightarrow} evaluate \ensuremath{\rightarrow} recover \ensuremath{\rightarrow} interpret}.
\item
  \textbf{A few datasets come back again and again.} Forgetting curves
  and signal detection theory (SDT) run through Weeks 2--13, so new
  methods are compared on familiar data. Chapter 6 adds categorization
  (GCM), response times (drift-diffusion) and reinforcement learning.
\item
  \textbf{Chapter weighting.} Ch1: 2 weeks · Ch2: 3 · Ch3: 2 · Ch4: 2 ·
  Ch5: 3 · Ch6: 2. One more week each goes to orientation, the midterm
  and project presentations.
\item
  \textbf{``Metropolis--Gibbs'' means two separate algorithms:}
  Metropolis--Hastings (Week 11) and Gibbs sampling (Week 12).
  MH-within-Gibbs is covered as an Advanced box.
\end{itemize}

\subsection{Weekly plan}

\begin{longtable}[]{@{}
  >{\raggedright\arraybackslash}p{(\linewidth - 12\tabcolsep) * \real{0.1429}}
  >{\raggedright\arraybackslash}p{(\linewidth - 12\tabcolsep) * \real{0.1429}}
  >{\raggedright\arraybackslash}p{(\linewidth - 12\tabcolsep) * \real{0.1429}}
  >{\raggedright\arraybackslash}p{(\linewidth - 12\tabcolsep) * \real{0.1429}}
  >{\raggedright\arraybackslash}p{(\linewidth - 12\tabcolsep) * \real{0.1429}}
  >{\raggedright\arraybackslash}p{(\linewidth - 12\tabcolsep) * \real{0.1429}}
  >{\raggedright\arraybackslash}p{(\linewidth - 12\tabcolsep) * \real{0.1429}}@{}}
\caption{\centering Weekly plan}\tabularnewline
\toprule\noalign{}
\begin{minipage}[b]{\linewidth}\raggedright
Wk
\end{minipage} & \begin{minipage}[b]{\linewidth}\raggedright
Ch
\end{minipage} & \begin{minipage}[b]{\linewidth}\raggedright
Lecture (1 h)
\end{minipage} & \begin{minipage}[b]{\linewidth}\raggedright
Lab (2 h)
\end{minipage} & \begin{minipage}[b]{\linewidth}\raggedright
Learning outcomes
\end{minipage} & \begin{minipage}[b]{\linewidth}\raggedright
Readings
\end{minipage} & \begin{minipage}[b]{\linewidth}\raggedright
Assignment / milestone
\end{minipage} \\
\midrule\noalign{}
\endfirsthead
\toprule\noalign{}
\begin{minipage}[b]{\linewidth}\raggedright
Wk
\end{minipage} & \begin{minipage}[b]{\linewidth}\raggedright
Ch
\end{minipage} & \begin{minipage}[b]{\linewidth}\raggedright
Lecture (1 h)
\end{minipage} & \begin{minipage}[b]{\linewidth}\raggedright
Lab (2 h)
\end{minipage} & \begin{minipage}[b]{\linewidth}\raggedright
Learning outcomes
\end{minipage} & \begin{minipage}[b]{\linewidth}\raggedright
Readings
\end{minipage} & \begin{minipage}[b]{\linewidth}\raggedright
Assignment / milestone
\end{minipage} \\
\midrule\noalign{}
\endhead
\bottomrule\noalign{}
\endlastfoot
1 & 1 & \textbf{Orientation: what is a cognitive model?} Why we model;
verbal vs.~formal theories; Marr's three levels; model families
(symbolic, connectionist, stochastic/dynamical, Bayesian); course
roadmap; \textbf{notation table} &
Simulate a random walk and compare its predictions with observations &
Distinguish descriptive, process and normative models; explain why
verbal theories are underspecified; run and plot a model simulation &
F\&L ch.~1--2; Verguts, introductory chapter {[}verify{]}; Anderson,
introductory chapter {[}verify{]} &
A0 (ungraded): one-page note: Pick a verbal theory; list what it leaves unspecified \\
2 & 1 & \textbf{Anatomy of a model:} data, free vs.~fixed parameters,
model equations, noise (data) model; simulation as a thinking tool;
first models: SDT (\(d'\), \(c\)) and exponential vs.~power forgetting &
Formulate and evaluate SDT (\(d'\), \(c\), ROC curves); simulate exponential and
power forgetting; explore how each parameter changes the predictions &
Identify the components of a formal model; compute \(d'\) and \(c\);
predict how parameters shape model output & F\&L ch.~2; L\&W ch.~11
(SDT, a preview) & \textbf{A1 released:} ``From verbal to formal'' (due
Wk 4) \\
3 & 2 & \textbf{Least squares \& optimization:} SSE, MSE, RMSD; error
surfaces; grid search, Nelder--Mead, gradient methods; local minima,
starting values, bounds and reparameterization &
Fit forgetting curves by grid search; plot the error surface; multi-start fitting
&
derive and evaluate an LS fit; compare optimizers; diagnose failed fits
& F\&L ch.~3; Ghodousian (direct search methods) & --- \\
4 & 2 & \textbf{Maximum likelihood:} probability vs.~likelihood;
log-likelihood; Bernoulli, binomial and Gaussian likelihoods; LS as MLE
under Gaussian noise & MLE for forgetting data (binomial recall counts)
and for SDT; compare MLE and LS estimates &
Derive a likelihood; analyze NLL minimization; explain when LS and MLE agree
& F\&L ch.~4; Azimifar \& Ahmadi (ML estimation) & \textbf{A1 due · A2
released:} LS vs.~MLE (due Wk 6) \\
5 & 2 & \textbf{Uncertainty \& multiple participants:} Hessian standard
errors, profile likelihood, parametric bootstrap; individual
vs.~aggregate fitting; aggregation artifacts & Bootstrap CIs; averaging
exponential learners produces a spurious ``power law of practice'';
individual fits & Quantify estimation uncertainty; explain aggregation
artifacts; choose a fitting level & F\&L ch.~5 & --- \\
6 & 3 & \textbf{Evaluation I, fit vs.~generalization:} \(R^2\), RMSD and
deviance and their limits; flexibility and overfitting; hold-out, k-fold
and LOO cross-validation & Polynomial overfitting demo; cross-validated
comparison of three forgetting models &
Explain why good fit \ensuremath{\neq} good model; analyze k-fold CV; interpret CV results
& F\&L ch.~10; Azimifar \& Ahmadi (generalization) & \textbf{A2 due · A3
released:} comparison + recovery (due Wk 9) \\
7 & 3 & \textbf{Evaluation II, selection \& recovery:} likelihood-ratio
test; AIC, AICc, BIC, Akaike weights; parameter and model recovery;
identifiability and trade-offs & Recovery study: true vs.~recovered
parameter plots; model-recovery confusion matrix (AIC/BIC) & Interpret
information criteria; design a recovery study; detect
non-identifiability & F\&L ch.~10 & \textbf{Project teams (1--2
students) formed; proposal brief released} \\
8 & 1--3 & \textbf{Midterm review:} integrated case study of the full
pipeline &
Midterm exam (lab slot, 90 min): conceptual questions + model interpretation, Ch. 1--3
& --- & Review of Weeks 1--7 & \textbf{Midterm} \\
9 & 4 & \textbf{Bayesian modeling I:} prior, likelihood, posterior,
evidence; beta-binomial conjugacy; grid approximation; posterior
summaries and HDI; credible vs.~confidence intervals & Grid posterior
for a recall probability; 2-D grid for SDT; vary prior and sample size;
prior predictive simulation & Compute a grid posterior; explain the role
of the prior; interpret credible intervals & K ch.~5--6; L\&W ch.~1, 3;
MKL ch.~1; F\&L ch.~6; Azimifar \& Ahmadi (Bayesian estimation) &
\textbf{A3 due · Project proposal due} \\
10 & 4 (+3) & \textbf{Bayesian modeling II:} posterior predictive
checks; marginal likelihood and \textbf{Bayes factors}; Savage--Dickey;
prior sensitivity; BIC as an approximation & PPCs for the forgetting
model; Savage--Dickey BF for \(d' > 0\); BF between retention models via
grid integration & Conduct a PPC; compute and interpret a Bayes factor;
evaluate prior sensitivity & L\&W ch.~7; F\&L ch.~11; K ch.~10 &
\textbf{A4 released:} Bayesian analysis by grid (due Wk 12) \\
11 & 5 & \textbf{Monte Carlo \& Metropolis--Hastings:} Monte Carlo
integration; Markov chains; MH; proposal tuning, acceptance rate,
burn-in & \textbf{MH from scratch} for the forgetting model; trace
plots; compare with the grid posterior &
Explain why MCMC works; analyze MH; tune a proposal & K ch.~7; F\&L
ch.~7 & --- \\
12 & 5 & \textbf{Gibbs sampling \& convergence diagnostics:} full
conditionals; Gibbs for a normal model; MH-within-Gibbs (Advanced);
multiple chains, \textbf{R-hat}, ESS, autocorrelation &
Gibbs from scratch ; analyze R-hat and ESS and check against ; break convergence and diagnose it
&
analyze Gibbs; interpret R-hat, ESS and trace plots; judge chain trustworthiness
& K ch.~7; MKL ch.~2; F\&L ch.~7 & \textbf{A4 due · A5 released:} MCMC +
diagnostics (due Wk 14) \\
13 & 5 &
Probabilistic \& hierarchical models: ; HMC/NUTS conceptually (Advanced); partial pooling and shrinkage (linked to Week 5
&
Hierarchical SDT in : diagnostics, PPCs, comparison with individual MLE fits; (Advanced) LOO/WAIC
&
Specify a model; explain partial pooling; report a full Bayesian workflow
& L\&W ch.~11; K ch.~9; F\&L ch.~9; MKL ch.~2, 10 & \textbf{Project
progress check} (feedback only) \\
14 & 6 & \textbf{Modeling in cognitive science I:} categorization
(prototype vs.~exemplar; \textbf{GCM}); choice/RT
(\textbf{drift-diffusion}, EZ-diffusion) & Fit the GCM by MLE to 5--4
category-structure data; EZ-diffusion on simulated RTs; interpret the
parameters & Fit an exemplar model; map DDM parameters to cognitive
constructs & L\&W ch.~17; F\&L ch.~14 {[}verify{]}; Verguts (relevant
chapters {[}verify{]}) & \textbf{A5 due} \\
15 & 6 & \textbf{Modeling in cognitive science II:} Rescorla--Wagner and
Q-learning with softmax; Bayesian \emph{models of cognition}
vs.~Bayesian \emph{data analysis}; connectionist models and links to AI
& Fit Q-learning (\(\alpha\), \(\beta\)) to simulated bandit data;
parameter recovery; (Advanced) hierarchical version in & Fit an RL
model; assess identifiability; contrast Bayesian cognition with Bayesian
statistics & Verguts (learning/RL chapters {[}verify{]}); F\&L ch.~13
{[}verify{]}; Anderson (relevant chapters {[}verify{]}) & Final project
work \\
16 & --- & \textbf{Final project presentations}; wrap-up on good
modeling practice & Presentations continue; peer feedback & Communicate
and critique a modeling study & --- &
Presentations · report due before the final exam \\
Exam period & 1--6 &
Final exam (written, cumulative, emphasis on Ch. 4--6; formula sheet provided; model interpretation, no laptops
& --- & --- & --- & \textbf{Final exam} \\
\end{longtable}

\textbf{Abbreviations:} F\&L = Farrell \& Lewandowsky (2018) · L\&W =
Lee \& Wagenmakers (2014) · K = Kruschke (2014) · MKL = Martin, Kumar \&
Lao (2021). Chapter numbers were written from memory. The ones marked
{[}verify{]} are less certain; please check all of them against the
books.

\subsection{Grading (0--20 scale)}

\begin{longtable}[]{@{}
  >{\raggedright\arraybackslash}p{(\linewidth - 2\tabcolsep) * \real{0.5000}}
  >{\raggedright\arraybackslash}p{(\linewidth - 2\tabcolsep) * \real{0.5000}}@{}}
\caption{\centering Grading (0--20 scale)}\tabularnewline
\toprule\noalign{}
\begin{minipage}[b]{\linewidth}\raggedright
Component
\end{minipage} & \begin{minipage}[b]{\linewidth}\raggedright
Points
\end{minipage} \\
\midrule\noalign{}
\endfirsthead
\toprule\noalign{}
\begin{minipage}[b]{\linewidth}\raggedright
Component
\end{minipage} & \begin{minipage}[b]{\linewidth}\raggedright
Points
\end{minipage} \\
\midrule\noalign{}
\endhead
\bottomrule\noalign{}
\endlastfoot
Lab check-offs (weekly in-lab exercises) & 2.0 \\
Assignments A1--A5 (1.0 each) & 5.0 \\
Short in-lecture quizzes (\ensuremath{\approx}10, drop the lowest 2) & 1.0 \\
Midterm exam (Week 8) & 3.0 \\
Final project: proposal 1.0 · presentation 2.0 · report + reproducible 3.0
& 6.0 \\
Final exam (cumulative) & 3.0 \\
\textbf{Total} & \textbf{20.0} \\
\end{longtable}

\subsection{Final project (overview; the full brief and rubric come in
Stage 3)}

\begin{itemize}
\tightlist
\item
  \textbf{Teams of 1--2.} The scope grows with team size. Everyone fits
  at least two candidate models, runs one model comparison and one
  parameter-recovery analysis. \textbf{Pairs also} do a Bayesian/MCMC
  analysis (or a hierarchical model) and submit an individual
  contribution statement.
\item
  \textbf{Data:} simulated data, an open dataset, or the student's own
  experimental data.
\end{itemize}
